\subsection{\texorpdfstring{Directed sets in \(L^{p}\) and increasing sequences in \(L^{p}\)}{Directed sets in L to the p and increasing sequences in L to the p}}

We have defined (§2, No.~6) an order relation \(\widetilde{f} \leqslant \widetilde{g}\) in the set \(\widetilde{\mathcal{F}}\) of equivalence classes of numerical functions defined and finite almost everywhere in \(X\) ; equipped with this order relation and with its vector space structure, \(\widetilde{\mathcal{F}}\) is a Riesz space. The corollary of Prop. 12 of No.~5 shows that if \(\widetilde{f}\) and \(\widetilde{g}\) are two elements of the subspace \(L to the p\) of \(\widetilde{\mathcal{F}}\) , then the supremum \(\sup (\widetilde{f}, \widetilde{g})\) of \(\widetilde{f}\) and \(\widetilde{g}\) in \(\widetilde{\mathcal{F}}\) (which is the class of each of the functions \(\sup (f, g)\) , where \(f \in \widetilde{f}\) and \(g \in \widetilde{g}\) ) belongs to \(L to the p\) ; this proves in particular that \(L to the p\) , equipped with the order relation induced by that of \(\widetilde{\mathcal{F}}\) , is a Riesz space.

PROPOSITION 13. --- In the Riesz space \(L^{p}\) , equipped with the topology defined by the norm \(\|f\|_{p}\) , the mapping \(\widetilde{f} \mapsto |f|\) is uniformly continuous, and the set of elements \(f \geqslant 0\) is closed.

The first part of the proposition follows at once from Prop. 11 of No.~5; since the set of \(\widetilde{f} \geqslant 0\) is also the set of \(\widetilde{f}\) such that \(|\widetilde{f}| = \widetilde{f}\) , it is closed, because \(\widetilde{f} \mapsto |\widetilde{f}|\) is a continuous mapping and \(\mathbf{L}^p\) is Hausdorff.

We thus see that the topology on \(L^{p}\) defined by the norm \(\|f\|_{p}\) is compatible with the ordered vector space structure of \(L^{p}\) (TVS, II, §2, No.~7).

PROPOSITION 14. --- Let H be a subset of the Riesz space \(L^{p}\) , consisting of classes \(\geqslant 0\) and directed for the relation \(\leqslant\) . For H to have a supremum in \(L^{p}\) , it is necessary and sufficient that

\[
\sup _ {\widetilde {f} \in \mathrm{H}} \| \widetilde {f} \| _ {p} <   + \infty .
\]

The supremum of H in \(L^{p}\) is then the limit (in the Banach space \(L^{p}\) ) of the section filter of H.

The condition is clearly necessary, since \(\widetilde{f} \mapsto \| \widetilde{f} \|_p\) is an increasing function on the set of elements \(\geqslant 0\) of \(L to the p\) . To see that it is sufficient, we first observe that it implies that the image of H under the mapping \(\widetilde{f} \mapsto \| \widetilde{f} \|_p\) has a limit in \(\mathbf{R}\) , by the monotone limit theorem; the image of the section filter \(\mathfrak{F}\) of H under this mapping is therefore a base of a Cauchy filter on \(\mathbf{R}\) . The proof will be complete if we show that \(\mathfrak{F}\) itself is a base of a Cauchy filter on \(L to the p\) ; for, \(\mathfrak{F}\) will then converge in \(L to the p\) , since \(L to the p\) is complete (No.~4, Th. 2), and the proposition will follow from TVS, II, §2, No.~7, Prop. 18.

To see that \(\mathfrak{F}\) is a base of a Cauchy filter, we shall make use of the following lemma:

Lemma. --- If \(f\) and \(g\) are two functions in \(\mathcal{L}^p\) such that \(0 \leqslant f \leqslant g\) , then

\[
\left(\mathrm{N} _ {p} (g - f)\right) ^ {p} \leqslant \left(\mathrm{N} _ {p} (g)\right) ^ {p} - \left(\mathrm{N} _ {p} (f)\right) ^ {p}.\tag{9}
\]

When \(f\) and \(g\) are continuous with compact support, the relation (9) may be written

\[
\int (g - f) ^ {p} d | \mu | \leqslant \int g ^ {p} d | \mu | - \int f ^ {p} d | \mu |
\]

and is then a consequence of the elementary inequality \((g - f)^{p} \leqslant g^{p} - f^{p}\) (No.~1, formula (2)). To pass from this to the general case, it suffices to observe that the two members of (9) are continuous functions on \(L^{p} \times L^{p}\) , and that every function \(f \geqslant 0\) in \(L^{p}\) is the limit (for convergence in mean of order p) of a sequence of continuous functions \(\geqslant 0\) with compact support, by the continuity of the mapping \(g \mapsto |g|\) on \(L^{p}\) (Prop. 11).

The lemma having been established, for every \(\varepsilon > 0\) there exists by hypothesis an \(\widetilde{f} \in \mathbf{H}\) such that, for every \(\widetilde{g} \geqslant \widetilde{f}\) belonging to \(\mathbf{H}\) , we have \((\|\widetilde{g}\|_p)^p - (\|\widetilde{f}\|_p)^p \leqslant \varepsilon\) ; from this it follows that \((\|\widetilde{g} - \widetilde{f}\|_p)^p \leqslant \varepsilon\) ; thus, if \(\widetilde{g}_1\) and \(\widetilde{g}_2\) are two elements in \(\mathbf{H}\) that are \(\geqslant \widetilde{f}\) , then \(\|\widetilde{g}_1 - \widetilde{g}_2\|_p \leqslant 2\varepsilon^{1/p}\) , which proves that \(\mathfrak{F}\) is a Cauchy filter base on \(\mathrm{L}^p\) and completes the proof of Prop. 14.

COROLLARY 1. --- If \(\widetilde{g}\) is the supremum of H in \(\mathbf{L}^p\) , then

\[
\| \widetilde {g} \| _ {p} = \lim _ {\widetilde {f} \in \mathrm{H}} \| \widetilde {f} \| _ {p} = \sup _ {\widetilde {f} \in \mathrm{H}} \| \widetilde {f} \| _ {p}.\tag{10}
\]

This follows from the continuity of the norm \(\|\widetilde{f}\|_{p}\) in \(L^{p}\) , and the monotone limit theorem.

COROLLARY 2. --- The Riesz space \(L^{p}\) is fully lattice-ordered.

Every directed set H in \(L^{p}\) (for the relation \(\leqslant\) ), consisting of classes \(\geqslant 0\) and bounded above in \(L^{p}\) , has a supremum: for, if \(\widetilde{h}\) is an upper bound for H in \(L^{p}\) , then \(\|f\|_{p} \leqslant \|h\|_{p}\) for all \(\widetilde{f} \in H\) , and Prop. 14 is applicable. This proves the corollary (Ch. II, §1, No.~3, Prop. 1).

The conclusions of Prop. 14 no longer hold when they are formulated for the functions in \(\mathcal{L}^p\) instead of their classes. To be precise, if M is a subset of \(\mathcal{L}^p\) , consisting of functions \(\geqslant 0\) , directed for the relation \(\leqslant\) , and such that \(\sup_{f\in\mathbb{M}}\mathrm{N}_{p}(f)<+\infty\) , the class of the upper envelope \(g\) of M is not

necessarily identical to the supremum in \(\mathrm{L}^p\) of the classes of the functions \(f \in \mathbf{M}\) ; in particular, \(g\) is not necessarily \(p\) -th power integrable, and even if \(g \in \mathcal{L}^p\) , \(\mathrm{N}_p(g)\) can be distinct from \(\sup_{f \in \mathbf{M}} \mathrm{N}_p(f)\) (cf.~§1, No.~3, Remark 1 following Th. 3).

Nevertheless, we have the following theorem:

THEOREM 5. --- Let \((f_n)\) be an increasing sequence of functions \(\geqslant 0\) in \(\mathcal{L}^p\) . For the upper envelope \(f\) of this sequence to be \(p\) -th power integrable, it is necessary and sufficient that \(\sup_{n} N sub p(f_n) < +\infty\) . The sequence \((f_n)\) is then convergent in mean of order \(p\) to \(f\) , and

\[
\mathrm{N} _ {p} (f) = \sup _ {n} \mathrm{N} _ {p} (f _ {n}) = \lim _ {n \rightarrow \infty} \mathrm{N} _ {p} (f _ {n}).\tag{11}
\]

The condition being obviously necessary, we need only prove that it is sufficient. Now, if the condition is satisfied then Prop. 14 shows that the sequence \((\tilde{f}_n)\) is a Cauchy sequence in \(L to the p\) , therefore the sequence \((f_n)\) is a Cauchy sequence in \(\mathcal{L}^p\) ; since \(f_n(x)\) tends to \(f(x)\) for all \(x \in X\) , \(f\) is \(p\) -th power integrable and is the limit of the sequence \((f_n)\) for the topology of convergence in mean of order \(p\) (No.~4, Cor. 1 of Th. 3). Therefore \(\mathrm{N}_p(f_n)\) tends to \(\mathrm{N}_p(f)\) since \(\mathrm{N}_p\) is a continuous function on \(\mathcal{L}^p\) .

COROLLARY 1. --- Let \((f_n)\) be a decreasing sequence of functions \(\geqslant 0\) in \(\mathcal{L}^p\) ; then, the lower envelope \(f\) of the sequence belongs to \(\mathcal{L}^p\) , the sequence \((f_n)\) converges in mean of order \(p\) to \(f\) , and

\[
\mathrm{N} _ {p} (f) = \lim _ {n \rightarrow \infty} \mathrm{N} _ {p} (f _ {n}) = \inf _ {n} \mathrm{N} _ {p} (f _ {n}).
\]

The first two assertions follow from Th. 5 applied to the sequence \(g_{n} = f_{1} - f_{n}\) , which is increasing and bounded above; the rest is then obvious.

COROLLARY 2. --- Let \((f_n)\) be a sequence of functions in \(\mathcal{L}^p\) . For the upper envelope \(f\) of the sequence \((f_n)\) to be \(p\) -th power integrable, it is necessary and sufficient that there exist a function \(g \geqslant 0\) such that \(\int^{*} g^p d|\mu| < +\infty\) and \(f_n \leqslant g\) for all \(n\) .

The condition is obviously necessary, on taking \(g = f^{+}\) . Conversely, suppose it is verified, and set \(g_{n} = \sup_{k\leqslant n}f_{k}\) ; the sequence \((g_{n})\) is increasing and consists of \(p\) -th power integrable functions (No.~5, Cor. of Prop. 12). The increasing sequence of positive functions \(h_n = g_n + g_1^-\) satisfies the conditions of Th. 5, since \(\mathrm{N}_p(h_n)\leqslant \mathrm{N}_p(g + g_1^-) < +\infty\) ; its upper envelope \(\sup_{n}h_{n}\) is therefore \(p\) -th power integrable, and the same is true of \(f = \sup_{n}h_{n} - g_{1}^{-}\) .

COROLLARY 3. --- Let A be a countable set, \(\mathfrak{F}\) a filter on A having a countable base, \((f_{\alpha})_{\alpha \in \mathrm{A}}\) a family of functions \(\geqslant 0\) in \(\mathcal{L}^p\) . Assume that there exists a function \(g \geqslant 0\) such that \(\mathrm{N}_p(g) < +\infty\) and \(f_{\alpha} \leqslant g\) for all \(\alpha \in \mathrm{A}\) ; then the function \(\limsup_{\mathfrak{F}} f_{\alpha}\) is \(p\) -th power integrable and

\[
\limsup _ {\mathfrak {F}} \mathrm{N} _ {p} (f _ {\alpha}) \leqslant \mathrm{N} _ {p} (\limsup _ {\mathfrak {F}} f _ {\alpha}).\tag{12}
\]

Let \((\mathbf{A}_n)\) be a decreasing base of \(\mathfrak{F}\) and set \(g_{n} = \sup_{\alpha \in \mathbf{A}_{n}}f_{\alpha}\) ; since \(\mathbf{A}_n\) is a countable set, it follows from Cor. 2 that \(g_{n}\) is \(p\) -th power integrable; on the other hand, \(\mathrm{N}_p(g_n)\geqslant \sup_{\alpha \in \mathbf{A}_n}\mathrm{N}_p(f_\alpha)\) . This being so, \(\limsup_{\mathfrak{F}}f_{\alpha}\) is the lower envelope of the decreasing sequence \((g_n)\) ; thus \(\limsup_{\mathfrak{F}}f_{\alpha}\) is \(p\) -th power integrable by Cor. 1, and

\[
\begin{array}{l} \mathrm{N} _ {p} (\limsup _ {\mathfrak {F}} f _ {\alpha}) = \mathrm{N} _ {p} \left(\inf _ {n} g _ {n}\right) = \lim _ {n \to \infty} \mathrm{N} _ {p} (g _ {n}) \\ \geqslant \lim _ {n \to \infty} \left(\sup _ {\alpha \in \mathrm{A} _ {n}} \mathrm{N} _ {p} (f _ {\alpha})\right) = \limsup _ {\mathfrak {F}} \mathrm{N} _ {p} (f _ {\alpha}). \end{array}
\]

